% GENERATED FILE. Edit canonical lessons, then run npm run generate:books.
% canonical-source-hash: fnv1a64:ec7cd3cb9b07fcbb
% canonical-lessons: TA-C537-cirpam, TA-C537-cetukku, TA-C537-uli, TA-C537-merukerru, TA-C537-cutu, TA-R537-review-530-537

\chapter{Clay and Stone}
\label{ch:ta-clay-and-stone}
\hlchaptermodality{537}

\begin{chapteropening}
\textbf{By the end of this chapter:} \emph{I can talk about carving, polishing and firing.}

Complete the last lesson of chapter 537: I can talk about carving, polishing and firing.
\end{chapteropening}

\section[\texorpdfstring{\ta{சிற்பம்} (ciṟpam)}{சிற்பம் (ciṟpam)}]{\ta{சிற்பம்} (ciṟpam) — a sculpture}

\label{lesson:TA-C537-cirpam}



\begin{quote}
Before the new one: say the Tamil for to paint a picture, then the Tamil for an exhibition.
\end{quote}

\subsection*{You'll want to know: \ta{சிற்பம்}}
\textbf{\ta{சிற்பம்}} — \emph{ciṟpam} — \textquotedblleft{}a sculpture\textquotedblright{}.

\begin{quote}
\ta{கோயிலின்} \ta{சுவர்களில்} \ta{பல} \ta{சிற்பங்கள்} \ta{இருக்கின்றன}.
\end{quote}

\begin{quote}
\emph{kōyiliṉ cuvarkaḷil pala ciṟpaṅkaḷ irukkiṉṟaṉa.}
\end{quote}

\begin{quote}
There are many sculptures on the walls of the temple.
\end{quote}

A sculptor is a \ta{சிற்பி}. The plural is \ta{சிற்பங்கள்}.

\begin{cousinweb}
\textbf{\ta{சிற்பம்}} (\emph{ciṟpam}) is borrowed from Sanskrit \emph{śilpa}, \textquotedblleft{}artistic work, a manual craft\textquotedblright{} (CDIAL 12470); Tamil wrote the \emph{lp} as \ta{ற்ப}.
\end{cousinweb}

\begin{grammarlens}[title={What you've built}]
One more word for describing what you lived and felt.
\end{grammarlens}

\subsection*{Your turn}
\begin{itemize}
\raggedright
  \item \emph{Say it:} \emph{ciṟpam}
  \item \emph{Say it:} \emph{ciṟpam}, once more
  \item \emph{Say it:} \emph{kōyiliṉ cuvarkaḷil pala ciṟpaṅkaḷ irukkiṉṟaṉa.} — There are many sculptures on the walls of the temple.
  \item \emph{Recall:} say the Tamil for to paint a picture, then the Tamil for an exhibition, then say \emph{ciṟpam} again
\end{itemize}

\begin{tcolorbox}[breakable,colback=teal!4,colframe=teal!35!black,title={Before you move on}]
What does \ta{சிற்பம்} mean? (\textquotedblleft{}A sculpture\textquotedblright{}.) Say it once more.
\end{tcolorbox}
\section[\texorpdfstring{\ta{செதுக்கு} (cetukku)}{செதுக்கு (cetukku)}]{\ta{செதுக்கு} (cetukku) — to carve}

\label{lesson:TA-C537-cetukku}



\begin{quote}
Before the new one: say the Tamil for an exhibition, then the Tamil for a sculpture.
\end{quote}

\subsection*{You'll want to know: \ta{செதுக்கு}}
\textbf{\ta{செதுக்கு}} — \emph{cetukku} — \textquotedblleft{}to carve\textquotedblright{}.

\begin{quote}
\ta{சிற்பி} \ta{கல்லில்} \ta{ஒரு} \ta{யானையைச்} \ta{செதுக்கினார்}.
\end{quote}

\begin{quote}
\emph{ciṟpi kallil oru yāṉaiyaic cetukkiṉār.}
\end{quote}

\begin{quote}
The sculptor carved an elephant in stone.
\end{quote}

Past (I …): \ta{செதுக்கினேன்}, \emph{cetukkiṉēṉ}.

The material takes -\ta{இல்}, the figure -\ta{ஐ}: \ta{கல்லில்} \ta{யானையைச்} \ta{செதுக்கு}.

\begin{cousinweb}
\textbf{\ta{செதுக்கு}} (\emph{cetukku}) is native Dravidian: DEDR 1953 gives \ta{செதுக்கு} (\emph{cetukku}), \textquotedblleft{}to cut off a surface, to pare, to hew with a tool\textquotedblright{}.
\end{cousinweb}

\begin{grammarlens}[title={What you've built}]
One more word for describing what you lived and felt.
\end{grammarlens}

\subsection*{Your turn}
\begin{itemize}
\raggedright
  \item \emph{Say it:} \emph{cetukku}
  \item \emph{Say it:} \emph{cetukku}, once more
  \item \emph{Say it:} \emph{ciṟpi kallil oru yāṉaiyaic cetukkiṉār.} — The sculptor carved an elephant in stone.
  \item \emph{Recall:} say the Tamil for an exhibition, then the Tamil for a sculpture, then say \emph{cetukku} again
\end{itemize}

\begin{tcolorbox}[breakable,colback=teal!4,colframe=teal!35!black,title={Before you move on}]
What does \ta{செதுக்கு} mean? (\textquotedblleft{}To carve\textquotedblright{}.) Say it once more.
\end{tcolorbox}
\section[\texorpdfstring{\ta{உளி} (uḷi)}{உளி (uḷi)}]{\ta{உளி} (uḷi) — a chisel}

\label{lesson:TA-C537-uli}



\begin{quote}
Before the new one: say the Tamil for a sculpture, then the Tamil for to carve.
\end{quote}

\subsection*{You'll want to know: \ta{உளி}}
\textbf{\ta{உளி}} — \emph{uḷi} — \textquotedblleft{}a chisel\textquotedblright{}.

\begin{quote}
\ta{தச்சர்} \ta{உளியால்} \ta{மரத்தைச்} \ta{செதுக்கினார்}.
\end{quote}

\begin{quote}
\emph{taccar uḷiyāl marattaic cetukkiṉār.}
\end{quote}

\begin{quote}
The carpenter carved the wood with a chisel.
\end{quote}

The tool takes -\ta{ஆல்}, \textquotedblleft{}with\textquotedblright{}: \ta{உளியால்}.

\begin{cousinweb}
\textbf{\ta{உளி}} (\emph{uḷi}) is native Dravidian: DEDR 699 gives \ta{உளி} (\emph{uḷi}), \textquotedblleft{}a chisel\textquotedblright{}, with cognates in Malayalam, Kota and Kannada.
\end{cousinweb}

\begin{grammarlens}[title={What you've built}]
One more word for describing what you lived and felt.
\end{grammarlens}

\subsection*{Your turn}
\begin{itemize}
\raggedright
  \item \emph{Say it:} \emph{uḷi}
  \item \emph{Say it:} \emph{uḷi}, once more
  \item \emph{Say it:} \emph{taccar uḷiyāl marattaic cetukkiṉār.} — The carpenter carved the wood with a chisel.
  \item \emph{Recall:} say the Tamil for a sculpture, then the Tamil for to carve, then say \emph{uḷi} again
\end{itemize}

\begin{tcolorbox}[breakable,colback=teal!4,colframe=teal!35!black,title={Before you move on}]
What does \ta{உளி} mean? (\textquotedblleft{}A chisel\textquotedblright{}.) Say it once more.
\end{tcolorbox}
\section[\texorpdfstring{\ta{மெருகேற்று} (merukēṟṟu)}{மெருகேற்று (merukēṟṟu)}]{\ta{மெருகேற்று} (merukēṟṟu) — to polish; to give a shine}

\label{lesson:TA-C537-merukerru}



\begin{quote}
Before the new one: say the Tamil for to carve, then the Tamil for a chisel.
\end{quote}

\subsection*{You'll want to know: \ta{மெருகேற்று}}
\textbf{\ta{மெருகேற்று}} — \emph{merukēṟṟu} — \textquotedblleft{}to polish; to give a shine\textquotedblright{}.

\begin{quote}
\ta{பொற்கொல்லர்} \ta{வளையலை} \ta{மெருகேற்றினார்}.
\end{quote}

\begin{quote}
\emph{poṟkollar vaḷaiyalai merukēṟṟiṉār.}
\end{quote}

\begin{quote}
The goldsmith polished the bangle.
\end{quote}

Past (I …): \ta{மெருகேற்றினேன்}, \emph{merukēṟṟiṉēṉ}.

Also of skills: \ta{திறமையை} \ta{மெருகேற்று}, to polish a skill.

\begin{cousinweb}
\textbf{\ta{மெருகேற்று}} (\emph{merukēṟṟu}) is a native compound: \ta{மெருகு} (\emph{meruku}), \textquotedblleft{}smoothness, lustre, polish\textquotedblright{} (DEDR 5074), and \ta{ஏற்று} (\emph{ēṟṟu}), \textquotedblleft{}to raise, to put on\textquotedblright{}.
\end{cousinweb}

\begin{grammarlens}[title={What you've built}]
One more word for describing what you lived and felt.
\end{grammarlens}

\subsection*{Your turn}
\begin{itemize}
\raggedright
  \item \emph{Say it:} \emph{merukēṟṟu}
  \item \emph{Say it:} \emph{merukēṟṟu}, once more
  \item \emph{Say it:} \emph{poṟkollar vaḷaiyalai merukēṟṟiṉār.} — The goldsmith polished the bangle.
  \item \emph{Recall:} say the Tamil for to carve, then the Tamil for a chisel, then say \emph{merukēṟṟu} again
\end{itemize}

\begin{tcolorbox}[breakable,colback=teal!4,colframe=teal!35!black,title={Before you move on}]
What does \ta{மெருகேற்று} mean? (\textquotedblleft{}To polish; to give a shine\textquotedblright{}.) Say it once more.
\end{tcolorbox}
\section[\texorpdfstring{\ta{சுடு} (cuṭu)}{சுடு (cuṭu)}]{\ta{சுடு} (cuṭu) — to bake; to fire (clay)}

\label{lesson:TA-C537-cutu}



\begin{quote}
Before the new one: say the Tamil for a chisel, then the Tamil for to polish.
\end{quote}

\subsection*{You'll want to know: \ta{சுடு}}
\textbf{\ta{சுடு}} — \emph{cuṭu} — \textquotedblleft{}to bake; to fire (clay)\textquotedblright{}.

\begin{quote}
\ta{குயவர்} \ta{பானைகளைச்} \ta{சூளையில்} \ta{சுட்டார்}.
\end{quote}

\begin{quote}
\emph{kuyavar pāṉaikaḷaic cūḷaiyil cuṭṭār.}
\end{quote}

\begin{quote}
The potter fired the pots in the kiln.
\end{quote}

Past (I …): \ta{சுட்டேன்}, \emph{cuṭṭēṉ}.

The past doubles the stop: \ta{சுட்டேன்}. The same verb means \textquotedblleft{}to be hot\textquotedblright{}, as in \ta{சூடு} (\emph{cūṭu}).

\begin{cousinweb}
\textbf{\ta{சுடு}} (\emph{cuṭu}) is native Dravidian: DEDR 2654 gives \ta{சுடு} (\emph{cuṭu}), \textquotedblleft{}to be hot, to burn, to roast, to bake, to burn bricks in a kiln\textquotedblright{}.
\end{cousinweb}

\begin{grammarlens}[title={What you've built}]
One more word for describing what you lived and felt.
\end{grammarlens}

\subsection*{Your turn}
\begin{itemize}
\raggedright
  \item \emph{Say it:} \emph{cuṭu}
  \item \emph{Say it:} \emph{cuṭu}, once more
  \item \emph{Say it:} \emph{kuyavar pāṉaikaḷaic cūḷaiyil cuṭṭār.} — The potter fired the pots in the kiln.
  \item \emph{Recall:} say the Tamil for a chisel, then the Tamil for to polish, then say \emph{cuṭu} again
\end{itemize}

\begin{tcolorbox}[breakable,colback=teal!4,colframe=teal!35!black,title={Before you move on}]
What does \ta{சுடு} mean? (\textquotedblleft{}To bake; to fire (clay)\textquotedblright{}.) Say it once more.
\end{tcolorbox}
\section[(dialogue)]{Looking back: from Children's Games to Clay and Stone}

\label{lesson:TA-R537-review-530-537}



\begin{quote}
Say the Tamil for hide and seek and a spinning top.
\end{quote}

\subsection*{Your turn}
Say these aloud:

\begin{itemize}
\raggedright
  \item hide and seek, a spinning top, marbles, carrom, a pinwheel
  \item basketball, badminton, volleyball, tennis, kabaddi
  \item boxing, wrestling, stick fencing, to kick, to hurl
  \item a race, to chase, to slide, to leap, to dive
  \item to beat a rival, a draw, a referee, a medal, a champion
  \item to take part, a pitch, a goal in sport, to take aim, a bat
  \item an art, a painting, a paintbrush, to paint a picture, an exhibition
  \item a sculpture, to carve, a chisel, to polish, to bake
\end{itemize}

\begin{tcolorbox}[breakable,colback=teal!4,colframe=teal!35!black,title={Before you move on}]
Hide and seek? (\textbf{\ta{கண்ணாமூச்சி}}, \emph{kaṇṇāmūcci}.) A spinning top? (\textbf{\ta{பம்பரம்}}, \emph{pamparam}.) Marbles? (\textbf{\ta{கோலி}}, \emph{kōli}.) Carrom? (\textbf{\ta{கேரம்}}, \emph{kēram}.) A pinwheel? (\textbf{\ta{காற்றாடி}}, \emph{kāṟṟāṭi}.) Basketball? (\textbf{\ta{கூடைப்பந்து}}, \emph{kūṭaippantu}.) Badminton? (\textbf{\ta{இறகுப்பந்து}}, \emph{iṟakuppantu}.) Volleyball? (\textbf{\ta{கைப்பந்து}}, \emph{kaippantu}.) Tennis? (\textbf{\ta{டென்னிஸ்}}, \emph{ṭeṉṉis}.) Kabaddi? (\textbf{\ta{கபடி}}, \emph{kapaṭi}.) Boxing? (\textbf{\ta{குத்துச்சண்டை}}, \emph{kuttuccaṇṭai}.) Wrestling? (\textbf{\ta{மல்யுத்தம்}}, \emph{malyuttam}.) Stick fencing? (\textbf{\ta{சிலம்பம்}}, \emph{cilampam}.) To kick? (\textbf{\ta{உதை}}, \emph{utai}.) To hurl? (\textbf{\ta{எறி}}, \emph{eṟi}.) A race? (\textbf{\ta{பந்தயம்}}, \emph{pantayam}.) To chase? (\textbf{\ta{துரத்து}}, \emph{turattu}.) To slide? (\textbf{\ta{சறுக்கு}}, \emph{caṟukku}.) To leap? (\textbf{\ta{தாவு}}, \emph{tāvu}.) To dive? (\textbf{\ta{முக்குளி}}, \emph{mukkuḷi}.) To beat a rival? (\textbf{\ta{தோற்கடி}}, \emph{tōṟkaṭi}.) A draw? (\textbf{\ta{சமநிலை}}, \emph{camanilai}.) A referee? (\textbf{\ta{நடுவர்}}, \emph{naṭuvar}.) A medal? (\textbf{\ta{பதக்கம்}}, \emph{patakkam}.) A champion? (\textbf{\ta{வாகையாளர்}}, \emph{vākaiyāḷar}.) To take part? (\textbf{\ta{பங்கெடு}}, \emph{paṅkeṭu}.) A pitch? (\textbf{\ta{ஆடுகளம்}}, \emph{āṭukaḷam}.) A goal in sport? (\textbf{\ta{கோல்}}, \emph{kōl}.) To take aim? (\textbf{\ta{குறிவை}}, \emph{kuṟivai}.) A bat? (\textbf{\ta{மட்டை}}, \emph{maṭṭai}.) An art? (\textbf{\ta{கலை}}, \emph{kalai}.) A painting? (\textbf{\ta{ஓவியம்}}, \emph{ōviyam}.) A paintbrush? (\textbf{\ta{தூரிகை}}, \emph{tūrikai}.) To paint a picture? (\textbf{\ta{தீட்டு}}, \emph{tīṭṭu}.) An exhibition? (\textbf{\ta{கண்காட்சி}}, \emph{kaṇkāṭci}.) A sculpture? (\textbf{\ta{சிற்பம்}}, \emph{ciṟpam}.) To carve? (\textbf{\ta{செதுக்கு}}, \emph{cetukku}.) A chisel? (\textbf{\ta{உளி}}, \emph{uḷi}.) To polish? (\textbf{\ta{மெருகேற்று}}, \emph{merukēṟṟu}.) To bake? (\textbf{\ta{சுடு}}, \emph{cuṭu}.)
\end{tcolorbox}
